\section{Discussion, Limitations, and Reproducibility Notes}

\subsection{Cross-Question Synthesis}

Across both questions, a consistent pattern emerges: discriminative linear models (Logistic Regression) provide stronger balanced performance when feature interactions matter. This was evident in sparse lexical email data and became even clearer in heterogeneous URL features.

However, the analysis also confirms that Naive Bayes remains a strong baseline with favorable simplicity-speed tradeoffs. It can produce highly precise alerts in some settings, which may be useful in staged triage pipelines.

\subsection{Engineering Fixes Required for Completion}

Completing this assignment to production-quality required several non-trivial fixes beyond filling placeholder cells:
\begin{enumerate}
    \item \textbf{Deterministic data loading:} all notebook loaders now use fixed project paths under \texttt{datasets/q1} and \texttt{datasets/q2}.
    \item \textbf{Label normalization guards:} explicit mappings for \texttt{spam/ham}, \texttt{Prediction}, and \texttt{phishing/legitimate} avoid silent type drift.
    \item \textbf{Safe stratification logic:} split strategy now handles edge cases cleanly instead of crashing.
    \item \textbf{Dependency handling:} runtime installation cells were replaced by deterministic import guards and environment-level dependency setup.
    \item \textbf{Figure export pipeline:} all report plots are saved with stable filenames for direct LaTeX inclusion.
\end{enumerate}

These changes converted the notebook from a partially interactive artifact to a reproducible experiment pipeline.

\subsection{Threats to Validity}

Despite successful completion, several limitations remain:
\begin{itemize}
    \item \textbf{Single split evaluation:} results are based on one train/test split; cross-validation would reduce variance.
    \item \textbf{Feature scope:} URL-only features exclude external signals such as DNS, WHOIS age, and content rendering behavior.
    \item \textbf{Model complexity:} only linear and Naive Bayes baselines were explored; tree-based ensembles or calibrated boosting could improve detection.
    \item \textbf{Cost-sensitive tuning:} thresholds were not optimized for domain-specific misclassification costs.
\end{itemize}

\subsection{Ablation-Oriented Interpretation}

Feature family ablation in Q2 provides a clear hierarchy:
\begin{enumerate}
    \item Combined features provide the best overall performance.
    \item Creative handcrafted features outperform purely structural or simple statistical signals.
    \item Statistical/content-only features are informative but insufficient alone for robust classification.
\end{enumerate}

This pattern supports the intuition that phishing URLs are best captured through multi-view representation rather than one-dimensional heuristics.

\subsection{Reproducibility Checklist}

The final artifact satisfies the following reproducibility criteria:
\begin{itemize}
    \item One-command notebook execution from project root in the provided virtual environment.
    \item Zero runtime cell errors in a full top-to-bottom run.
    \item Deterministic generation of all expected report figures.
    \item IEEE report compilation from local LaTeX sources without manual post-processing.
\end{itemize}

These properties ensure that another evaluator can re-run the entire pipeline and recover both notebook outputs and PDF report figures.
